\chapter{Glosario}

\begin{description}
    \item[Acuerdo de Conclusión:] Documento final que cierra el expediente, pudiendo resultar en una recomendación, un cese de actos o un pronunciamiento de improcedencia.
    \item[Acuse de Recibo:] Comprobante digital generado por el sistema que confirma que la DDP ha recibido la solicitud, otorgando un folio único de seguimiento.
     \item[Análisis morfosintáctico:] Proceso de PLN que analiza la estructura gramatical
    y función sintáctica de cada palabra en una oración para extraer información
    relevante del texto.
    \item[Arquitectura de Cero Conocimiento (Zero-Knowledge):] Diseño de seguridad donde el servidor no tiene acceso a las llaves de cifrado del usuario, siendo incapaz de ver los datos almacenados.
    \item[Arquitectura de Microservicios:] Estilo arquitectónico que estructura una aplicación como un conjunto de servicios pequeños, autónomos y acoplados de forma mínima.
    \item[Arquitectura Monolítica:] Modelo de software donde todos los componentes del sistema están entrelazados en un solo programa, dificultando su escalabilidad.
    \item[AUC-ROC:] Curva que mide la capacidad del modelo para distinguir entre clases; un área cercana a 1 indica un desempeño óptimo en la clasificación.
    \item[Bolsa de Palabras (BoW):] Técnica de representación de texto basada en la frecuencia de aparición de cada palabra, sin considerar el orden gramatical.
    \item[Comunidad Politécnica:] Conjunto de estudiantes, docentes, personal administrativo y egresados del IPN.
    \item[Competencia:] Facultad o atribución legal que tiene un órgano (DDP) para conocer y resolver un asunto específico basado en su normativa.
    \item[DDP (Defensoría de los Derechos Politécnicos):] Órgano encargado de velar por los derechos de la comunidad del IPN.
     \item[Dicotómico:] Que admite únicamente dos valores mutuamente excluyentes.
    \item[Disponibilidad:] Capacidad de un sistema de permanecer operativo y accesible durante un periodo de tiempo determinado.
     \item[Embeddings:] Representaciones numéricas de palabras o textos en un espacio
    vectorial donde la proximidad geométrica refleja similitud semántica.
    \item[Escalabilidad:] Propiedad de un sistema para manejar una cantidad creciente de trabajo o su potencial para ser ampliado.
    \item[Estado del Arte:] Análisis de los conocimientos, tecnologías y soluciones más avanzados que existen actualmente sobre un tema de investigación.
    \item[Expediente Digital:] Repositorio electrónico que centraliza documentación, evidencias multimedia y actuaciones jurídicas de un caso.
    \item[F1-Score:] Métrica reina que representa el promedio balanceado entre Precisión y Recall, útil para evaluar la calidad global del modelo.
    \item[Flujo Operativo:] Secuencia estandarizada de etapas por las que debe pasar una queja, desde su recepción hasta su archivo definitivo.
    \item[Geolocalización:] Técnica que permite asignar una ubicación geográfica exacta (coordenadas) a un reporte o evento.
    \item[GRC (Governance, Risk, and Compliance):] Estrategia para gestionar el gobierno corporativo, los riesgos y el cumplimiento legal de una organización.
    \item[Hiperplano:] Frontera o muro matemático en un espacio multidimensional que separa las distintas clases (ej. quejas competentes de improcedentes).
    \item[Integridad Documental:] Garantía de que los documentos que integran el expediente digital no han sido manipulados desde su carga original.
    \item[IPN (Instituto Politécnico Nacional):] Institución pública mexicana de educación media superior, superior y posgrado.
    \item[Java Spring Boot:] Framework diseñado para simplificar el desarrollo de aplicaciones Java basadas en microservicios.
    \item[Ley Orgánica:] Marco legal máximo que rige la estructura, fines y organización del Instituto Politécnico Nacional.
    \item[Manual de Organización / Procedimientos:] Documentos normativos que definen las funciones y los pasos operativos internos de la Defensoría.
    \item[Margen:] Distancia entre los vectores de soporte y el hiperplano; el sistema busca maximizar esta distancia para reducir errores de clasificación.
    \item[Matriz de Confusión:] Tabla que permite visualizar el desempeño de un algoritmo, mostrando los aciertos y las confusiones (falsos positivos y negativos).
    \item[PLN / NLP (Procesamiento de Lenguaje Natural):] Rama de la IA que permite a las máquinas entender, interpretar y manipular el lenguaje humano.
     \item[Pooling:] Operación que condensa una secuencia de vectores en una
    representación única de dimensión fija.
    \item[PostgreSQL:] Sistema de gestión de bases de datos relacional orientado a objetos y de código abierto.
    \item[Precisión:] Métrica que indica qué tan exacto es el sistema cuando clasifica una queja como positiva, evitando falsas alarmas.
    \item[Primer Contacto:] Etapa inicial donde se recibe la solicitud y se realiza el filtrado preliminar de la queja.
    \item[Quejoso:] Persona que interpone una queja o denuncia por una presunta vulneración de sus derechos.
    \item[Radicación:] Acto formal de registrar y asignar un número de expediente a una queja para iniciar su proceso legal.
    \item[RBAC (Role-Based Access Control):] Método de control de acceso que restringe o permite funciones del sistema según el rol asignado al usuario.
    \item[Recall (Sensibilidad):] Capacidad del sistema para detectar todos los casos positivos reales, asegurando que ninguna queja válida sea ignorada.
    \item[Reglas de Negocio:] Políticas institucionales que el software debe ejecutar automáticamente para garantizar la validez del proceso.
    \item[Similitud de Coseno:] Medida métrica utilizada para determinar qué tan parecidos son dos documentos de texto en un espacio vectorial.
    \item[Soberanía de Datos:] Principio que establece que los datos están sujetos a las leyes y control de la institución donde se generan.
    \item[Stemming y Lematización:] Procesos lingüísticos para reducir palabras a su raíz o forma base (lema) para facilitar el análisis.
    \item[SVM (Support Vector Machine):] Algoritmo de aprendizaje supervisado que busca el hiperplano óptimo para clasificar datos en diferentes categorías.
    \item[Términos Procesales / Plazos Legales:] Lapsos de tiempo obligatorios otorgados por la normativa para realizar una acción jurídica.
    \item[TF-IDF:] Técnica de pesaje que resalta palabras clave de un problema y resta importancia a palabras comunes o irrelevantes.
    \item[Tokenización:] Proceso de fragmentar una cadena de texto en unidades más pequeñas llamadas tokens (palabras o símbolos).
    \item[Trazabilidad:] Capacidad de seguir el historial, la aplicación y la ubicación de un expediente a través de todas sus etapas.
    \item[Truco del Kernel (Kernel Trick):] Técnica que proyecta datos a dimensiones superiores para encontrar una separación lineal donde antes no era posible.
    \item[Validación Cruzada (K-Fold):] Método de evaluación que divide los datos en $K$ partes para asegurar que el sistema sea robusto con diferentes subconjuntos de datos.
    \item[Variables de Holgura ($\xi$):] Parámetro de flexibilidad en un modelo SVM que permite que algunos puntos crucen la frontera para manejar datos con ruido.
    \item[Vectores de Soporte:] Puntos o quejas que se encuentran más cerca de la frontera de decisión y definen la posición del hiperplano.
    \item[Whistleblowing:] Acto de denunciar internamente o ante las autoridades actividades ilegales o poco éticas dentro de una organización.
    \item[Word Embeddings:] Representación de palabras como vectores en un espacio multidimensional donde la cercanía geométrica indica cercanía semántica.
    \item[Word2Vec / GloVe:] Modelos específicos utilizados para generar embeddings basados en el contexto local o estadísticas globales del texto.
\end{description}